\subsection{Concentration du travail entre rangs}\label{sec:concentration}

Le maximum de charge décrit un seul rang. Pour savoir combien de rangs portent le calcul, on classe leurs charges par ordre décroissant, \(W_{(1)}\ge\cdots\ge W_{(48)}\). La part des cellules attribuées aux \(k\) premiers rangs est
\[
  C(k)=\frac{\sum_{r=1}^{k}W_{(r)}}{\sum_{r=1}^{48}W_{(r)}}.
\]
Un partage égal suit la droite \(C(k)=k/48\). La figure~\ref{fig:rank-work} compare trois politiques sur les mêmes 8\,128 paires, en ordre de longueur des fragments. Chaque paire contribue le produit des deux longueurs, quelle que soit la durée de son alignement.

\chart{distributed/concentration_rangs.png}{Part des cellules calculées par les rangs les plus chargés sur deux nœuds. Le ruban donne l’intervalle de confiance ponctuel à 95\,\% sur cinq lancements avec vol. La barre à douze rangs donne l’écart-type. Les répartitions statiques se reproduisent à l’identique.}{fig:rank-work}{python/studies/rank_concentration.py}
Avec le vol, les douze rangs les plus chargés exécutent en moyenne \(56{,}5\,\%\) des cellules (écart-type \(2{,}0\) points, IC à 95\,\% de \(54{,}0\) à \(59{,}0\,\%\)). La répartition varie selon les transferts de paires entre rangs. La section suivante conserve les politiques d’attribution et change le moteur de calcul à l’intérieur de chaque rang.

\synthese{Sur 48 rangs, les douze premiers reçoivent \(77{,}0\,\%\) du travail avec des blocs de Hilbert de même effectif. Pondérer ces blocs par le nombre de cellules ramène leur part à \(25{,}2\,\%\), pour un quart des rangs.}
